\documentclass{amsart}
\usepackage{amsmath,amssymb,amsthm,hyperref}
\usepackage{mathtools}

\newtheorem{theorem}{Theorem}
\newtheorem{lemma}{Lemma}
\newtheorem{proposition}{Proposition}
\newtheorem{corollary}{Corollary}
\theoremstyle{definition}
\newtheorem{definition}{Definition}
\newtheorem{remark}{Remark}

\newcommand{\E}{\mathbb{E}}
\newcommand{\R}{\mathbb{R}}
\newcommand{\kB}{k_{B}}
\newcommand{\Wext}{\langle W_{\mathrm{ext}}\rangle}
\newcommand{\dd}{\mathrm{d}}
\DeclareMathOperator{\KL}{D}

\begin{document}

\title{Stochastic thermodynamics of feedback flashing ratchets:
information bound, maximal entropy reduction, and tight work/power/efficiency bounds}
\author{}
\date{}
\maketitle

\section{Set-up and standing assumptions}

We make the dynamical model precise so that every later statement is
unambiguous.

\paragraph{Particle dynamics.}
A Brownian particle has position $X_t\in\R$ obeying the overdamped Langevin
equation
\begin{equation}\label{eq:langevin}
\gamma\,\dd X_t = -\,a_t\,V'(X_t)\,\dd t + \sqrt{2\gamma \kB T}\;\dd B_t ,
\end{equation}
where $\gamma>0$ is the friction coefficient, $T>0$ the bath temperature,
$B_t$ a standard Wiener process, and
\[
V:\R\to\R,\qquad V(x+L)=V(x),\qquad V\in C^{1},
\]
a fixed $L$-periodic asymmetric potential. The \emph{control variable}
\[
a_t\in\{0,1\}
\]
is piecewise constant in $t$: $a_t=1$ means the potential is ON,
$a_t=0$ means it is OFF (free diffusion). Equation \eqref{eq:langevin} is
understood in the It\^o sense; existence and uniqueness of a strong solution
on each interval where $a_t$ is constant follow from the Lipschitz/linear-growth
property of $V'$ implied by $V\in C^1$ and periodicity (so $V'$ is bounded).

\paragraph{Feedback protocol.}
Let $[0,\tau]$ be the operating window. A \emph{feedback protocol}
$\pi$ is a rule that produces the control trajectory
$a_{[0,\tau]}:=(a_t)_{t\in[0,\tau]}$ as a (possibly randomised) functional
of measurements of the particle. We treat the two standard cases jointly.

\begin{itemize}
\item \emph{Discrete-time feedback.} At sampling instants
$0\le t_1<t_2<\dots<t_N\le\tau$ the controller obtains measurement outcomes
$Y_k=m(X_{t_k},\xi_k)$, where $m$ is a measurement channel and $\xi_k$ is
internal measurement noise (so $Y_k$ may be noisy/coarse-grained). After the
$k$-th measurement the controller chooses the control value on
$[t_k,t_{k+1})$ as a measurable, possibly randomised function of the
\emph{record} $Y^{(k)}:=(Y_1,\dots,Y_k)$. Thus $a_t$ on $[t_k,t_{k+1})$ is a
function of $Y^{(k)}$.

\item \emph{Continuous-time feedback.} The controller has access to the
filtration $\mathcal F^Y_t=\sigma(Y_s:s\le t)$ generated by a (possibly noisy)
observation process $Y_s$, and $a_t$ is $\mathcal F^Y_t$-adapted (causal,
nonanticipating).
\end{itemize}

In both cases the control sequence is \emph{non-Markovian} in the sense that
$a_t$ generally depends on the whole past record, not only on the current
state.

\paragraph{Thermodynamic quantities.}
For a control trajectory $a_{[0,\tau]}$ the stochastic work injected by the
external agent that flashes the potential is, by the first law for the
flashing ratchet (the potential changes only at the switching instants, where
the particle position is frozen),
\begin{equation}\label{eq:work}
W \;=\; \sum_{\text{switch times }s} \bigl(a_{s^+}-a_{s^-}\bigr)\,V(X_s)
\;=\; \int_0^\tau V(X_t)\,\dd a_t ,
\end{equation}
the last expression being a Stieltjes integral against the pure-jump process
$a_t$. The work \emph{extracted} is $W_{\mathrm{ext}}=-W$. The free energy of
the ON state at canonical equilibrium is
$F=-\kB T\log\!\int_0^L e^{-V(x)/\kB T}\dd x$ (per period), and $\Delta F$
denotes the change of free energy between the initial and final ensembles.

Throughout, $\langle\cdot\rangle$ denotes expectation over both the thermal
noise $B$ and the measurement/protocol randomness.

\section{(a) The information quantity for an arbitrary protocol}

\subsection{Reference (open-loop) dynamics and the path measure}
Fix the operating window and an initial law $\rho_0$ for $X_0$. A control
trajectory $a_{[0,\tau]}$ together with the Langevin dynamics
\eqref{eq:langevin} defines a probability measure on particle paths. The
crucial object is the joint law of \emph{(particle path, control record)}.

Let $\Omega$ be the path space of $X$, let $\mathcal Y$ be the space of
measurement records (the vector $Y^{(N)}$ in discrete time, or the path
$Y_{[0,\tau]}$ in continuous time), and let
\[
P(\dd x,\dd y)
\]
be the joint law generated by the actual feedback protocol $\pi$.

We compare $\pi$ to a \emph{reference protocol} that produces an identically
distributed control trajectory but \emph{without using the measurement to
decide it}: i.e.\ the control is drawn from the same marginal law of
$a_{[0,\tau]}$, but \emph{independently} of $X$. Concretely, let
$P_X(\dd x)$ be the marginal law of the particle path and $P_Y(\dd y)$ the
marginal law of the control record under $\pi$. The reference (``no
feedback / blind replay'') joint law is the product
\[
P^{\mathrm{ref}}(\dd x,\dd y) := P_X(\dd x)\,P_Y(\dd y).
\]

\subsection{Definition of $I$}

\begin{definition}[Information of a feedback protocol]\label{def:I}
For a feedback protocol $\pi$ on the flashing ratchet over $[0,\tau]$, the
\emph{information} is the mutual information between the particle path and the
controller's record,
\begin{equation}\label{eq:Idef}
I \;:=\; \mathcal I\!\bigl(X_{[0,\tau]};\,Y\bigr)
\;=\; \KL\!\bigl(P \,\big\|\, P^{\mathrm{ref}}\bigr)
\;=\; \E\!\left[\log\frac{\dd P}{\dd P^{\mathrm{ref}}}(X,Y)\right]\ \ge 0,
\end{equation}
where $\KL(\cdot\|\cdot)$ is relative entropy and the equality with the
Kullback--Leibler divergence to the product of marginals is the defining
identity of mutual information.
\end{definition}

For \emph{discrete-time} protocols, $Y=Y^{(N)}=(Y_1,\dots,Y_N)$ and the chain
rule for mutual information together with causality of the control gives the
operationally transparent decomposition
\begin{equation}\label{eq:Ichain}
I=\sum_{k=1}^{N}\mathcal I\!\bigl(X_{[0,\tau]};\,Y_k \,\big|\, Y^{(k-1)}\bigr).
\end{equation}
This is exactly the \emph{transfer-entropy-type} sum appropriate to a
non-Markovian sequence of control actions: at the $k$-th step the controller
extracts $\mathcal I(X;Y_k\mid Y^{(k-1)})$ nats of \emph{new} information about
the trajectory, conditioned on everything learned so far. For
\emph{continuous-time} protocols, \eqref{eq:Idef} is the continuous-time
limit of \eqref{eq:Ichain}, i.e.\ a transfer-entropy \emph{rate} integrated
over $[0,\tau]$; when the observation is the full position with measurement
noise it is the directed-information / transfer entropy from $X$ to the
control.

\begin{remark}
$I$ depends only on the joint law of (trajectory, record) and is
manifestly nonnegative, finite when the measurement has finitely many
outcomes per step, and reduces to the ordinary Shannon mutual information for
a single measurement. It is the quantity appearing in the generalised second
law: for any feedback protocol on \eqref{eq:langevin},
\begin{equation}\label{eq:2ndlaw}
\Wext \;\le\; -\Delta F + \kB T\, I .
\end{equation}
We do not re-derive \eqref{eq:2ndlaw} (it is the Sagawa--Ueda generalised
second law, valid for any measurement-and-feedback protocol on a Langevin
system); we take it as the framework hypothesis stated in the problem and
\emph{compute the ratchet-specific value of $I$ and its maximum}. The
saturating condition of \eqref{eq:2ndlaw} is reversibility of the controlled
dynamics; we use it in \S4 to identify saturating protocols.
\end{remark}

\section{(b) Maximal entropy reduction $I_{\max}$}

The entropy production of the controlled system is bounded below by $-I$, so
the maximal \emph{entropy reduction} achievable by feedback equals the maximal
value of $I$ over admissible protocols. The decisive structural fact about the
flashing ratchet is that the controller's \emph{action space is binary}:
$a_t\in\{0,1\}$. We exploit this to bound $I$ by the entropy of the control
process, which is the only channel through which information feeds back into
the work \eqref{eq:work}.

\subsection{Reduction to the control variable}

\begin{lemma}[Data-processing reduction]\label{lem:dpi}
Let $A:=a_{[0,\tau]}$ denote the control trajectory generated by the protocol.
Then
\begin{equation}\label{eq:IleIA}
I \;=\; \mathcal I\!\bigl(X_{[0,\tau]};\,Y\bigr)\;\ge\;
\mathcal I\!\bigl(X_{[0,\tau]};\,A\bigr),
\end{equation}
and the work \eqref{eq:work}, hence $\Wext$, is a functional of $(X,A)$ only.
Consequently the thermodynamically \emph{usable} information is
$I_A:=\mathcal I(X_{[0,\tau]};A)$, and
\begin{equation}\label{eq:IAbound}
I_A \;\le\; H(A),
\end{equation}
where $H(A)$ is the Shannon entropy of the control trajectory.
\end{lemma}

\begin{proof}
The control trajectory $A$ is, by definition of the protocol, a
(measurable, possibly randomised) function of the record $Y$ alone:
$A=g(Y,\zeta)$ with $\zeta$ a randomisation variable independent of $X$. Hence
the Markov chain
\[
X \longrightarrow Y \longrightarrow A
\]
holds (given $Y$, $A$ is conditionally independent of $X$). The data-processing
inequality for mutual information gives
$\mathcal I(X;A)\le \mathcal I(X;Y)=I$, which is \eqref{eq:IleIA}.

That $\Wext$ depends on $(X,A)$ only is immediate from \eqref{eq:work}: the
integrand $V(X_t)$ and the integrator $\dd a_t$ involve only $X$ and $A$. The
measurement record $Y$ influences the work \emph{exclusively} through the
control $A$ it produces; this is the precise sense in which $I_A$ is the
usable part.

Finally, $\mathcal I(X;A)=H(A)-H(A\mid X)\le H(A)$ because conditional entropy
is nonnegative; this is \eqref{eq:IAbound}.
\end{proof}

\eqref{eq:IleIA}--\eqref{eq:IAbound} say that, although the raw measurement may
carry arbitrarily much information $I$, only the part transduced into the
control can be converted to work, and that part is capped by the entropy of the
binary switching process. Combining with \eqref{eq:2ndlaw}:
\begin{equation}\label{eq:2ndlaw2}
\Wext \;\le\; -\Delta F + \kB T\, I_A \;\le\; -\Delta F + \kB T\, H(A).
\end{equation}

\subsection{Bounding the entropy of the control trajectory}

We now bound $H(A)$ in the two protocol classes.

\begin{theorem}[Maximal entropy reduction, discrete time]\label{thm:Imax-discrete}
Suppose the controller makes $N$ measurements and may switch the potential at
most once after each measurement (one binary decision per measurement, the
standard flashing-ratchet protocol). Then for every feedback protocol
\begin{equation}\label{eq:Imax-discrete}
I_A \;\le\; H(A) \;\le\; N\log 2 \;=:\; I_{\max}^{(N)} .
\end{equation}
The bound is tight: it is attained, $H(A)=N\log 2$, by any protocol whose
binary decisions $D_k\in\{0,1\}$ (``set $a=$ ON or OFF on the $k$-th
interval'') are i.i.d.\ uniform; and the \emph{usable} bound $I_A=N\log2$ is
attained, in the idealised noiseless-measurement limit, by a protocol whose
$k$-th decision is a balanced ($\Pr=\tfrac12$ each) deterministic function of
$X_{t_k}$, e.g.\ a sign rule based on the local force.
\end{theorem}

\begin{proof}
The control trajectory is determined by the sequence of binary decisions
$D=(D_1,\dots,D_N)\in\{0,1\}^N$ (the value of $a$ on each inter-measurement
interval), because $a_t$ is piecewise constant and equals $D_k$ on
$[t_k,t_{k+1})$. Hence $A$ is a deterministic bijective recoding of $D$, so
$H(A)=H(D)$. By subadditivity of Shannon entropy and the fact that each $D_k$
is a binary random variable,
\[
H(D)\le\sum_{k=1}^N H(D_k)\le \sum_{k=1}^N \log 2 = N\log 2,
\]
using $H(D_k)\le\log 2$ with equality iff $D_k$ is uniform. With
\eqref{eq:IAbound} this proves \eqref{eq:Imax-discrete}.

\emph{Tightness of $H(A)=N\log2$:} if the $D_k$ are i.i.d.\ uniform on
$\{0,1\}$ then $H(D)=\sum_k H(D_k)=N\log2$, so the chain of inequalities is an
equality.

\emph{Tightness of $I_A=N\log 2$:} take $D_k=\phi(X_{t_k})$ for a measurable
$\phi:\R\to\{0,1\}$ chosen so that, under the (closed-loop) law of $X_{t_k}$,
$\Pr[\phi(X_{t_k})=1]=\tfrac12$ and the decisions are conditionally
independent given the relevant past (achievable, e.g., by re-randomising the
threshold so that the marginal of each $D_k$ is exactly balanced). With
noiseless measurement, $Y_k=X_{t_k}$ and $D_k$ is a deterministic function of
$X$, so $H(D\mid X)=0$ and
$I_A=\mathcal I(X;A)=H(A)-H(A\mid X)=H(D)=N\log2$. Thus the entropy reduction
bound is saturated. (With noisy measurement $D_k=\phi(Y_k)$ one has
$H(A\mid X)>0$ and the usable information is strictly below $N\log2$, by the
strict data-processing gap.)
\end{proof}

\begin{theorem}[Maximal entropy-reduction \emph{rate}, continuous time]
\label{thm:Imax-cont}
Let the continuous-time controller switch the potential a random number $K$ of
times on $[0,\tau]$, with a minimal dwell time $\delta>0$ between consecutive
switches (any physical actuator has a finite switching speed). Then the
entropy of the control trajectory, and hence the usable information, obeys
\begin{equation}\label{eq:Icont2}
I_A\;\le\;H(A)\;\le\; \log 2 \;+\; H(K)\;+\;
\E\!\Bigl[\log\binom{\lfloor\tau/\delta\rfloor}{K}\Bigr]
\;\le\;\log 2 + H(K) + \E[K]\,\log\tfrac{\tau}{\delta}.
\end{equation}
If in addition the expected number of switches is budgeted by $\E[K]\le R\tau$
for a fixed rate $R<\infty$ (unbounded switching costs unbounded actuation
work), then
\begin{equation}\label{eq:Icont}
I_A\;\le\;H(A)\;\le\; \log2 + H(K) + R\tau\,\log\tfrac{\tau}{\delta},
\end{equation}
so the entropy-reduction \emph{rate} is finite,
$\dot I_{\max}:=\limsup_{\tau\to\infty}H(A)/\tau \le R\,\log(\tau/\delta)$ for
fixed $\tau/\delta$ regime, and in any case $H(A)/\tau$ is bounded for finite
$R,\delta$. In the un-budgeted, zero-dwell idealisation ($\delta\to0$) the
switch-time entropy diverges and $I_A$ is unbounded, reflecting that an
idealised continuous noiseless controller can in principle extract unbounded
information; any finite actuation constraint renders it finite.
\end{theorem}

\begin{proof}
A right-continuous piecewise-constant trajectory $a_{[0,\tau]}$ with values in
$\{0,1\}$ is determined by (i) its initial value $a_0\in\{0,1\}$, (ii) the
number of switches $K$, and (iii) the ordered switch times
$0<s_1<\dots<s_K<\tau$. Indeed, given $a_0$ and the ordered switch times, the
value on each interval is determined by alternation, so the map
$(a_0,K,\{s_i\})\mapsto A$ is injective; therefore
\[
H(A)\le H(a_0,K,\{s_i\})\le H(a_0)+H(K)+H(\{s_i\}\mid K),
\]
by subadditivity and the chain rule of Shannon entropy, and no further entropy
is added by the (deterministic) interval values. The initial value contributes
$H(a_0)\le\log2$. With minimal dwell time $\delta$, every admissible switch
time lies on a grid of at most $\lfloor\tau/\delta\rfloor$ slots; given $K$,
the unordered set of switch times is one of at most
$\binom{\lfloor\tau/\delta\rfloor}{K}$ choices, so
\[
H(\{s_i\}\mid K)\le \E\!\Bigl[\log\binom{\lfloor\tau/\delta\rfloor}{K}\Bigr]
\le \E[K]\,\log\tfrac{\tau}{\delta},
\]
using $\binom{n}{k}\le n^{k}$ with $n=\lfloor\tau/\delta\rfloor\le\tau/\delta$.
This gives \eqref{eq:Icont2}. Imposing $\E[K]\le R\tau$ gives \eqref{eq:Icont}.
The data-processing inequality \eqref{eq:IAbound} gives $I_A\le H(A)$
throughout. Letting $\delta\to0$ makes $\log(\tau/\delta)\to\infty$, so the
switch-time entropy and hence the bound diverge, establishing the stated
idealised unboundedness.
\end{proof}

\begin{remark}[Interpretation of $I_{\max}$]
Theorems \ref{thm:Imax-discrete}--\ref{thm:Imax-cont} give the
ratchet-specific maximal entropy reduction: \emph{per binary control decision
the feedback can reduce the system entropy by at most $\log 2$} (one nat-valued
bit; restoring units, $\kB\log2$ of entropy). In discrete time with $N$
decisions, $I_{\max}=N\log2$; in continuous time with finite switching rate
$R$ and dwell time $\delta$, the entropy reduction stays finite and only
diverges in the unphysical zero-dwell limit. The binary nature of ``flashing''
is what caps the per-action information at one bit.
\end{remark}

\section{(c) Complete thermodynamic bounds: work, power, efficiency}

We now assemble the consequences. Let $A$ be the control process, $I_A$ its
usable information \eqref{eq:IleIA}, and recall \eqref{eq:2ndlaw2}.

\subsection{Extracted work}

\begin{theorem}[Tight work bound]\label{thm:work}
For every feedback protocol on the flashing ratchet,
\begin{equation}\label{eq:workbound}
\Wext \;\le\; -\Delta F + \kB T\,I_A \;\le\; -\Delta F + \kB T\,H(A),
\end{equation}
and in the cyclic/steady operation of a ratchet, where the potential and the
particle distribution return statistically to their starting condition over a
period so that $\Delta F=0$, this reads
\begin{equation}\label{eq:workbound-cyclic}
\Wext \;\le\; \kB T\, I_A \;\le\; \kB T\, H(A).
\end{equation}
In discrete time with $N$ binary decisions, $\Wext\le N\,\kB T\log2$. Equality
in the first inequality of \eqref{eq:workbound} holds iff the feedback-controlled
forward process is statistically reversible (the saturation condition of the
generalised second law); equality in the second holds iff $H(A\mid X)=0$, i.e.\
the control is a deterministic function of the trajectory (noiseless,
non-randomised feedback).
\end{theorem}

\begin{proof}
\eqref{eq:workbound} is \eqref{eq:2ndlaw2}; \eqref{eq:workbound-cyclic} is its
specialisation to $\Delta F=0$. The numerical bound follows from
Theorem~\ref{thm:Imax-discrete}. The first equality condition is the
equality condition of the Sagawa--Ueda second law \eqref{eq:2ndlaw}, namely
vanishing of the total (system $+$ controller) entropy production, i.e.\
reversibility of the controlled dynamics. The second equality condition is the
equality condition of \eqref{eq:IAbound}: $I_A=H(A)-H(A\mid X)=H(A)$ iff
$H(A\mid X)=0$, i.e.\ $A$ is $\sigma(X)$-measurable.
\end{proof}

\subsection{Output power}

The output power is the mean extracted work per unit time. For cyclic
operation with period $\tau$ and $N$ measurements per period:

\begin{corollary}[Power bound]\label{cor:power}
\begin{equation}\label{eq:power}
\mathcal P \;:=\; \frac{\Wext}{\tau}\;\le\;\frac{\kB T\,I_A}{\tau}
\;\le\;\frac{\kB T\,N\log2}{\tau}.
\end{equation}
Writing $f=N/\tau$ for the measurement frequency, the power is bounded by
$\kB T\,(\log2)\,f$, i.e.\ by the information \emph{rate} times $\kB T$. In
continuous time, by Theorem~\ref{thm:Imax-cont},
$\mathcal P\le \kB T\,H(A)/\tau$, which stays finite for any finite switching
rate $R$ and dwell time $\delta$.
\end{corollary}

\begin{proof}
Divide \eqref{eq:workbound-cyclic} by $\tau$ and insert the relevant
$I_{\max}$. The continuous-time statement divides $H(A)$ from
Theorem~\ref{thm:Imax-cont} by $\tau$.
\end{proof}

\subsection{Efficiency}

For an information-driven ratchet the natural thermodynamic efficiency
compares the work extracted to the information ``fuel'' supplied, both in
energy units. Define the \emph{information efficiency}
\begin{equation}\label{eq:eta-def}
\eta_{\mathrm{info}} \;:=\; \frac{\Wext + \Delta F}{\kB T\, I_A}\, .
\end{equation}

\begin{corollary}[Efficiency bound]\label{cor:eff}
For every feedback protocol with $I_A>0$,
\begin{equation}\label{eq:eta}
\eta_{\mathrm{info}} \;\le\; 1,
\end{equation}
with equality iff the controlled forward process is reversible (the saturation
condition of \eqref{eq:2ndlaw}). If, in addition, one charges the controller
the raw acquired information $I\ge I_A$ rather than the usable part $I_A$, the
realised efficiency
$\tilde\eta:=(\Wext+\Delta F)/(\kB T\,I)\le I_A/I\le1$ is further degraded by
the information-transduction loss $I_A/I$, which equals $1$ iff the control
loses no information about $X$ (lossless transduction, $X\to Y\to A$ with no
data-processing gap).
\end{corollary}

\begin{proof}
By \eqref{eq:workbound}, $\Wext+\Delta F\le \kB T I_A$; dividing by
$\kB T I_A>0$ gives \eqref{eq:eta}. Equality is the equality case of
\eqref{eq:2ndlaw}. The refined statement uses $\Wext+\Delta F\le\kB T I_A\le
\kB T I$ and divides by $\kB T I$; the factor $I_A/I\le1$ is exactly the
data-processing ratio from Lemma~\ref{lem:dpi}, equal to $1$ iff
$\mathcal I(X;A)=\mathcal I(X;Y)$, i.e.\ no information about $X$ is lost in
producing the control.
\end{proof}

\subsection{Saturating protocols}

\begin{proposition}[Identification of optimal protocols]\label{prop:opt}
A discrete-time feedback flashing ratchet simultaneously saturates the
work bound \eqref{eq:workbound}, the power bound \eqref{eq:power}, and the
efficiency bound \eqref{eq:eta} if and only if it satisfies all three of:
\begin{enumerate}
\item[\textup{(i)}] \emph{Maximal information acquisition}: each binary decision is
balanced and deterministic in the (noiseless) measurement,
$\Pr[D_k=1]=\tfrac12$, $D_k=\phi(X_{t_k})$, giving $I_A=N\log2$
(Theorem~\ref{thm:Imax-discrete});
\item[\textup{(ii)}] \emph{Lossless transduction}: $X\to Y\to A$ has no
data-processing gap, $\mathcal I(X;A)=\mathcal I(X;Y)$, i.e.\ the controller
discards no relevant information (automatic in the noiseless case
$Y_k=X_{t_k}$, $D_k=\phi(Y_k)$);
\item[\textup{(iii)}] \emph{Thermodynamic reversibility}: the controlled forward
process is statistically reversible, so the generalised second law
\eqref{eq:2ndlaw} holds with equality.
\end{enumerate}
\end{proposition}

\begin{proof}
By Theorem~\ref{thm:work}, the second inequality in \eqref{eq:workbound}
($I_A=H(A)=N\log2$) is equality iff (i); the first inequality is equality iff
(iii). The efficiency bound \eqref{eq:eta} is equality iff (iii). The
refined-efficiency / no-loss requirement is (ii) (Corollary~\ref{cor:eff}).
The power bound \eqref{eq:power} chains the same two inequalities divided by
$\tau$, so it is equality under (i)+(iii). Conversely each named condition is
exactly an equality case of one of the two inequalities, so all are necessary.
\end{proof}

\begin{remark}[Physical reading]
Conditions (i)--(iii) make the optimal feedback flashing ratchet completely
explicit: \emph{measure as informatively as the binary action allows (one
balanced bit per decision), feed all of it back without loss, and switch
quasi-statically/reversibly}. Under these conditions every period extracts
exactly $\kB T\log2$ of work per balanced bit, the power is $\kB T\log2$ times
the measurement rate, and the information efficiency is unity. The binary
``flashing'' constraint caps the per-action benefit at one bit
($\kB T\log2$ of free energy / work), which is the Landauer/Szilard value;
the maximal entropy reduction of the device is therefore $I_{\max}=N\log2$
(discrete) and finite for any physical actuation budget (continuous).
\end{remark}

\section*{Summary}
\begin{itemize}
\item[(a)] The information of a feedback flashing-ratchet protocol is the
mutual information $I=\mathcal I(X_{[0,\tau]};Y)$ between trajectory and
record, Definition~\ref{def:I}, with the causal/non-Markovian decomposition
\eqref{eq:Ichain} (a transfer-entropy sum; a transfer-entropy rate in
continuous time). The thermodynamically usable part is
$I_A=\mathcal I(X_{[0,\tau]};A)\le I$ (Lemma~\ref{lem:dpi}).
\item[(b)] The maximal entropy reduction is set by the binary control:
$I_{\max}=N\log2$ in discrete time (Theorem~\ref{thm:Imax-discrete}), and
remains finite for any finite switching rate and dwell time in continuous time
(Theorem~\ref{thm:Imax-cont}); both are saturated by balanced, deterministic,
noiseless feedback.
\item[(c)] Hence the tight bounds
$\Wext\le -\Delta F+\kB T I_A\le -\Delta F+\kB T N\log2$
(Theorem~\ref{thm:work}),
$\mathcal P\le \kB T I_A/\tau\le \kB T(\log2)N/\tau$
(Corollary~\ref{cor:power}), and $\eta_{\mathrm{info}}\le1$
(Corollary~\ref{cor:eff}), all simultaneously saturated exactly by the
protocols of Proposition~\ref{prop:opt}.
\end{itemize}

\end{document}
